% =============================================================================
%  PARTIE V — Des lumières pilotées par Ableton
% =============================================================================

\chapter{La refonte multi-projecteurs}
\label{ch:multi}

Le contrôleur d'origine ne pilotait qu'un seul projecteur. Pour produire de
vrais jeux de lumière, \Icare{} gère désormais un \textbf{patch} complet : une
liste de projecteurs, chacun décrit par un \emph{profil} (plan de canaux) et une
\emph{adresse} de départ dans l'univers DMX.

\section{Les profils de projecteurs}

Un profil dit quel canal fait quoi. \Icare{} embarqué les profils suivants,
qui correspondent au matériel :

\begin{table}[H]
  \centering
  \begin{tabular}{@{}llp{6.8cm}@{}}
    \toprule
    \textbf{Profil} & \textbf{Canaux} & \textbf{Détail} \\
    \midrule
    \code{rgbwa\_par\_5ch} & 5 & Rouge, Vert, Bleu, Blanc, Ambre. \\
    \addlinespace
    \code{wild\_wash\_132\_6ch} & 6 & Intensité, Strobe, Rouge, Vert, Bleu,
      sensibilité son. \\
    \addlinespace
    \code{cameo\_qspot40\_9ch} & 9 & Intensité, intensité fine, Strobe, R, G, B,
      Blanc, macro, température. \\
    \addlinespace
    \code{stairville\_mfog1500\_3ch} & 3 & Volume brouillard, intervalle, durée. \\
    \bottomrule
  \end{tabular}
  \caption{Profils livrés.}
\end{table}

\section{Le patch par défaut}

Par défaut, les appareils sont posés à la suite à partir du canal 1 :

\begin{table}[H]
  \centering
  \begin{tabular}{@{}llcc@{}}
    \toprule
    \textbf{Appareil} & \textbf{Profil} & \textbf{Canal de départ} & \textbf{Canaux} \\
    \midrule
    ParLED1 \dots ParLED6 & RGBWA 5 canaux & 1, 6, 11, 16, 21, 26 & 5 \\
    Wash1, Wash2 & Wild Wash 132 & 31, 37 & 6 \\
    QSpot1, QSpot2 & Cameo Q-Spot 40 & 43, 52 & 9 \\
    Fog & M-Fog 1500 & 61 & 3 \\
    \bottomrule
  \end{tabular}
  \caption{Patch par défaut (63 canaux utilisés).}
  \label{tab:patch}
\end{table}

Le patch est validé au chargement : un profil inconnu, un chevauchement de
canaux ou une adresse hors de l'univers sont refusés. Il vit dans
\file{patch.json} à côté de l'exécutable et peut être modifié depuis l'API
\api{POST /api/patch}.

\section{Comment le rendu gère les différences}

Chaque mode \og{}génératif\fg{} (couleur fixe, effets, aléatoire, audio)
produit une couleur et une intensité globales. Le moteur l'écrit sur chaque
projecteur :

\begin{itemize}[leftmargin=1.4em]
  \item si le profil possède un \textbf{canal d'intensité}, il est écrit tel
        quel et les couleurs en brut ;
  \item sinon (PAR RGBWA 5 canaux), les couleurs sont \textbf{multipliées par
        l'intensité} : c'est l'intensité qui pilote la luminosité.
\end{itemize}

Le brouillard possède ses propres canaux, indépendants des modes : on peut
l'activer séparément via \api{POST /api/fog}.

\begin{note}
Le Q-Spot 40 et le Wild Wash sont des projecteurs \textbf{fixes} : leurs profils
n'ont ni pan ni tilt. Les curseurs de position de l'onglet Direct ne servent
donc qu'à la visualisation 3D.
\end{note}

\chapter{Les « looks »}
\label{ch:looks}

Un \textbf{look} est un état de lumière statique que l'on peut rappeler : une
liste de cibles et, pour chacune, une couleur et une intensité.

\section{Cibles}

Une cible peut être :

\begin{itemize}[leftmargin=1.4em]
  \item \code{"*"} : tous les projecteurs ;
  \item le nom d'un \textbf{groupe} (\code{pars}, \code{wash}, \code{spots},
        \code{fx}) ;
  \item le nom d'un \textbf{projecteur} précis (\code{QSpot1}, \code{Fog}\dots).
\end{itemize}

Les cibles sont fusionnées dans l'ordre \code{"*"}, puis le groupe, puis le
nom : on met donc une base globale et on surcharge au besoin.

\section{Résolution}

\code{resolve()} transforme le look en valeurs de canaux concrètes. Pour un
projecteur sans canal d'intensité, l'intensité demandée \og{}multiplie\fg{} les
couleurs. Exemple de look :

\begin{lstlisting}[language=json,caption={Un look : base bleue, projecteurs rouges.}]
{
  "name": "Refrain",
  "targets": {
    "pars":  { "r": 0, "g": 0, "b": 255, "dimmer": 220 },
    "QSpot1": { "r": 255, "g": 80, "b": 0 }
  },
  "fog": { "level": 120 }
}
\end{lstlisting}

Les looks sont de simples fichiers JSON dans \file{Looks/}, créés depuis
l'onglet \textbf{Show} : on choisit un nom, une cible, une couleur et une
intensité, puis on sauvegarde.

\chapter{Le lien avec Ableton Live}
\label{ch:ableton}

\section{Principe}

\Icare{} écoute Ableton Live par le réseau grâce au script distant gratuit
\href{https://github.com/ideoforms/AbletonOSC}{AbletonOSC}. Aucun câble
supplémentaire : tout passe en OSC sur \code{127.0.0.1}.

\begin{figure}[H]
  \centering
  \begin{tikzpicture}[
      node distance=10mm,
      font=\small,
      box/.style={draw=icodeln,fill=icodebg,rounded corners,inner sep=7pt,align=center},
      flow/.style={-Latex,draw=icareblue,thick},
    ]
    \node[box] (abl) {Ableton Live\\\scriptsize + remote script AbletonOSC};
    \node[box,right=2.4cm of abl] (icare) {\Icare{} (Python)\\\scriptsize moteur + moteur de scènes};
    \node[box,right=2.4cm of icare] (dmx) {Enttec USB Pro\\\scriptsize $\rightarrow$ projecteurs};
    \draw[flow] (abl) -- node[above,draw=none,fill=none,font=\scriptsize]{OSC :11000} node[below,draw=none,fill=none,font=\scriptsize]{tempo, scène, beat} (icare);
    \draw[flow] (icare) -- node[above,draw=none,fill=none,font=\scriptsize]{DMX512} (dmx);
    \draw[flow,dashed] (icare.north) to[out=120,in=60] node[above,draw=none,fill=none,font=\scriptsize]{réponses :11001} (abl.north);
  \end{tikzpicture}
  \caption{Boucle Ableton $\rightarrow$ \Icare{} $\rightarrow$ DMX.}
\end{figure}

Ce qu'\Icare{} observe en continu :

\begin{table}[H]
  \centering
  \begin{tabular}{@{}ll@{}}
    \toprule
    \textbf{Adresse OSC} & \textbf{Information} \\
    \midrule
    \api{/live/song/get/tempo} & Tempo (BPM). \\
    \api{/live/song/get/is\_playing} & Lecture en cours. \\
    \api{/live/view/get/selected\_scene} & Scène sélectionnée. \\
    \api{/live/scene/get/name} & Nom de chaque scène. \\
    \api{/live/song/get/beat} & Compteur de temps (beat events). \\
    \bottomrule
  \end{tabular}
  \caption{Les principales données suivies.}
\end{table}

\section{Installation}

\begin{enumerate}[leftmargin=1.6em]
  \item Exécuter \file{tools\textbackslash install\_abletonosc.ps1} : le script
        télécharge AbletonOSC et le copie dans
        \file{Documents\textbackslash Ableton\textbackslash User Library\textbackslash Remote Scripts}.
  \item Redémarrer Live.
  \item \emph{Préférences $\rightarrow$ Link, Tempo \& MIDI} $\rightarrow$
        Control Surface = \textbf{AbletonOSC}.
  \item Dans \Icare{}, onglet \textbf{Show} $\rightarrow$ \textbf{Connecter
        Ableton}. Le témoin passe au vert.
\end{enumerate}

\begin{note}
Sans Ableton, le lien reste simplement \og{}déconnecté\fg{} : \Icare{} continue
de fonctionner normalement.
\end{note}

\section{Mapper les scènes aux looks}

L'onglet \textbf{Show} liste les scènes d'Ableton (récupérées via OSC). Pour
chacune, on choisit le look à déclencher. Dès qu'une scène est sélectionnée ou
lancée dans Ableton, \Icare{} applique le look correspondant.

\begin{lstlisting}[caption={Mapping persistant (scenes.json).},numbers=none]
{ "enabled": true,
  "mapping": { "INTRO": "intro_bleu", "A": "couplet", "CHORUS": "refrain" } }
\end{lstlisting}

\section{Produire son show}

Une méthode simple, adaptée au template Ableton du studio :

\begin{enumerate}[leftmargin=1.6em]
  \item Dans \Icare{}, créer un look par ambiance : \code{INTRO},
        \code{A}, \code{B}, \code{CHORUS}, \code{BREAK}, \code{SOLO},
        \code{OUTRO}.
  \item Mapper chaque scène Ableton au look du même nom, une fois pour toutes.
  \item En répétition ou en concert, il suffit de lancer les scènes dans
        Ableton : les lumières suivent.
  \item Réglage fin : le \emph{beat} d'Ableton pourra servir plus tard à caler
        les effets (strobe sur le temps, chenillard à la mesure).
\end{enumerate}

\chapter{Architecture logicielle mise à jour}
\label{ch:archi2}

\begin{table}[H]
  \centering
  \begin{tabular}{@{}lp{9cm}@{}}
    \toprule
    \textbf{Module} & \textbf{Rôle} \\
    \midrule
    \file{webapp/fixtures.py} & Profils de projecteurs et patch. \\
    \file{webapp/looks.py} & Modèle de look, résolution, stockage. \\
    \file{webapp/osc.py} & Encodeur/décodeur OSC minimal. \\
    \file{webapp/ableton.py} & Client AbletonOSC + moteur de scènes. \\
    \file{webapp/dmx\_controller.py} & Moteur multi-projecteurs (thread de sortie). \\
    \file{webapp/api\_show.py} & API patch / looks / scènes / Ableton. \\
    \file{webapp/templates/show.html} & Onglet Show. \\
    \bottomrule
  \end{tabular}
  \caption{Nouveaux modules de la refonte.}
\end{table}

Le moteur DMX garde la même surface publique qu'avant (\code{set\_values},
\code{strobe}, \code{effect\_*,} \code{audio\_push}, \code{health}, \dots), mais
les applique désormais au \textbf{patch entier}. Les anciennes fonctions de
l'onglet Direct et le lecteur de séquences continuent donc de fonctionner.

\begin{note}
Le contrôleur par défaut écrivait une seule adresse de couleur. Il écrit
maintenant une \emph{frame} complète : pour chaque projecteur, les canaux de son
profil sont remplis, ce qui rend possible un rendu différent d'un appareil à
l'autre.
\end{note}
